% =====================================================================
% Section 2: Related Work
% =====================================================================
%
% Status: Draft §2 Related Work v2 (2026-05-16)
%   v1: three-paragraph structure, 564 words, 18 citations
%   v2: per GPT review —
%        - HGRAG表述纠错(entity hypergraph,not community structure)
%        - propagation 描述去过度概括
%        - triples → facts(术语跟 §4 对齐;PropRAG 处保留)
%        - 删第一段末 "two failure modes" 提前披露
%        - monoT5 泛化为通用 neural reranker 表述
%        - 补 §2.1 各 GraphRAG 系统差异 + §2.3 coverage/diversity 关系
%        - 删 2 处冒号
%        - PropRAG cite key 改 wang2025proprag(待 bib finalize 核对)
%
% Structure: three-paragraph form (PropRAG/NeocorRAG/Relink-style)
%   §2.1 Graph-based Retrieval for Multi-hop QA
%   §2.2 KG Reasoning and Query-time Chain Mining
%   §2.3 Personalized PageRank and Coverage-aware Reranking
%
% Each paragraph ends with an "In contrast / \methodname{} differs ..."
% sentence positioning our differences without overclaiming prior work
% weaknesses.
%
% Dependencies (BibTeX) — needs finalization:
%   already in main paper:
%     angeli2015leveraging, gutierrez2025hipporag, sun2024tog,
%     huang2026relink, peng2026neocorrag
%   to add:
%     jimenez2024hipporag (HippoRAG)
%     wang2025proprag (PropRAG; verify author + year)
%     edge2024graphrag (Microsoft GraphRAG)
%     hgrag202x (HGRAG; verify exact venue/year)
%     haveliwala2002topic, jeh2003scaling
%     carbonell1998mmr, lin2011submodular
%     trivedi2023ircot, press2023selfask
%     karpukhin2020dpr, robertson2009bm25
%     baek2023kaping
%
% =====================================================================

\section{Related Work}
\label{sec:related}

\paragraph{Graph-based retrieval for multi-hop QA.}
Standard retrieval-augmented generation methods rank passages by dense
or lexical similarity to the question
\citep{karpukhin2020dpr,robertson2009bm25}, which is effective for
single-hop questions but tends to miss bridge passages whose surface
similarity to the question is low. Recent work augments this pipeline
with explicit graph structure to improve multi-hop coverage.
HippoRAG \citep{jimenez2024hipporag} and HippoRAG 2
\citep{gutierrez2025hipporag} extract OpenIE facts from the corpus and
run personalized PageRank over a corpus-wide graph of passage and
entity nodes seeded from the question. PropRAG \citep{wang2025proprag}
replaces triples with propositions and discovers multi-hop evidence by
beam search over a proposition graph, motivated by the observation that
triples can lose context. GraphRAG \citep{edge2024graphrag} builds
community-summary graphs over the corpus and supports retrieval through
hierarchical summary aggregation, while hypergraph-based methods such
as HGRAG \citep{hgrag2024} model entities and passages within a
hypergraph and propagate relevance across entity-level and
passage-level representations. These methods mainly differ in the graph
representation and in the search operator used over the pre-built
structure, ranging from corpus-wide PageRank to beam search over
propositions and hierarchical summary aggregation. \methodname{} adopts
the same offline OpenIE-based indexing strategy
\citep{angeli2015leveraging,gutierrez2025hipporag}, but performs
propagation on a query-local fact subgraph and then applies
relation-grounded expansion and coverage-aware reranking on top of the
local ranking.

\paragraph{KG reasoning and query-time chain mining.}
A second line of work uses extracted relations as a reasoning substrate
that is traversed or updated at query time. ToG \citep{sun2024tog}
performs LLM-guided traversal over a pre-built knowledge graph,
selecting paths step by step, while related KG-RAG approaches
\citep{baek2023kaping} retrieve subgraphs and pass them to a language
model as context. Relink \citep{huang2026relink} addresses KG
incompleteness by dynamically reconstructing query-specific evidence
graphs and repairing missing facts as needed. NeocorRAG
\citep{peng2026neocorrag} keeps documents as the reader unit but
invokes a language model on retrieved text to mine entity-relation
chains as prompt augmentation. Interleaved chain-of-thought retrieval
methods such as IRCoT \citep{trivedi2023ircot} and Self-Ask
\citep{press2023selfask} similarly drive retrieval through
language-model-generated reasoning traces, alternating between
generation and lookup. In contrast, \methodname{} keeps the OpenIE
structure on the indexing side. The offline graph is neither traversed
by an LLM at query time nor regenerated per question, and chain-aware
retrieval is computed by graph propagation rather than by query-time
chain generation over retrieved text. The reader consumes
natural-language passages, and any query-time language model use is
restricted to extracting question-side anchors and requirements.

\paragraph{Personalized PageRank and coverage-aware reranking.}
\methodname{}'s retrieval and enhancement components draw on
long-standing techniques in information retrieval. Personalized
PageRank \citep{haveliwala2002topic,jeh2003scaling} provides a
principled way to bias graph propagation toward a query, and has been
used in graph-augmented retrieval to score nodes given a seed
distribution. On the reranking side, MMR \citep{carbonell1998mmr}
introduced diversity-aware reordering through marginal relevance against
an already-selected set, and submodular set functions
\citep{lin2011submodular} formalize coverage and diminishing returns
when selecting summary or evidence sets; both perspectives motivate the
saturation behavior we adopt for $\mathrm{cov}_q$. Pointwise and
pairwise neural rerankers provide strong passage-level relevance
estimates, but they usually score candidates independently from the
coverage of the selected evidence set, and so do not directly model
which entities or relations remain uncovered. \methodname{} adapts both
the graph-propagation and the coverage-aware perspectives to the
multi-hop retrieval setting. Propagation is restricted to a query-induced
subgraph and runs over OpenIE fact units rather than over passage or
word nodes directly, while the coverage signal is computed under a
noisy-OR model over question-side entity-relation requirements rather
than over query tokens or word-level features. These choices align the
retrieval and reranking objectives with the question's relational
structure rather than with surface similarity alone.
